\documentclass[10pt,a4paper]{amsart}
\usepackage[margin=19mm]{geometry}
\usepackage{amsmath,amssymb,mathtools}
\usepackage[scr=rsfs,cal=euler]{mathalfa}
\usepackage{xfrac}
\usepackage[unicode,hidelinks,bookmarks=false]{hyperref}
\newcommand{\ie}{i.e.\ }
\newcommand{\cf}{cf.\ }
\newcommand{\resp}{respectively\ }
\newcommand{\Subsection}{\subsection}
\providecommand{\nobreakdash}{\nobreak\hbox{-}}
\setlength{\emergencystretch}{2em}
\multlinegap0pt
\usepackage{mdframed}
\usepackage{enumerate}
\usepackage{dsfont}
\usepackage{mathtools}
\usepackage{amssymb}
\usepackage{bm}
\usepackage{algorithm}\usepackage{algpseudocode}
\usepackage{xcolor}
\usepackage{tikz}
\newtheorem{prop}{Proposition}
\newtheorem{lem}{Lemma}
\newcommand{\Ber}{\operatorname{Ber}}
\def\E{\mathbb{E}}
\newcommand{\GKBM}{\operatorname{GKBM}}
\newcommand{\La}{\Lambda}
\def\P{\mathbb{P}}
\newcommand{\Poi}{\operatorname{Poi}}
\def\R{\mathbb{R}}
\newcommand{\al}{\alpha}
\def\bA{\bm{A}}
\newcommand{\be}{\beta}
\newcommand{\bsigma}{\bm{\sigma}}
\def\cE{\mathcal{E}}
\def\cL{\mathcal{L}}
\def\cV{\mathcal{V}}
\newcommand{\de}{\delta}
\newcommand{\eps}{\epsilon}
\newcommand{\ga}{\gamma}
\newcommand{\ka}{\kappa}
\newcommand{\la}{\lambda}
\def\weq{ \ = \ }
\def\wge{ \ \ge \ }
\def\wle{ \ \le \ }
\newcommand{\zeronorm}[1]{{\lVert#1\rVert}_0}
\title{Community Detection on Block Models with Geometric Kernels}
\author{Konstantin Avrachenkov, B. R. Vinay Kumar, Lasse Leskel\"a}
\date{Source: arXiv:2403.02802v2 (CC BY 4.0)}
\begin{document}
\maketitle

\noindent\emph{Source and licence.} Community Detection on Block Models with Geometric Kernels. Version arXiv:2403.02802v2 (CC BY 4.0), distributed by arXiv under the Creative Commons Attribution 4.0 International licence, http://creativecommons.org/licenses/by/4.0/, as recorded in the arXiv licence metadata for this exact version and shown on the abstract page https://arxiv.org/abs/2403.02802v2. Full licence notice: \texttt{licenses/arxiv-2403.02802-CC-BY-4.0.txt}.\par\smallskip

\noindent\textit{Editorial scope.} Counted unit: the article's prop prop:origin\_bad (source0.tex lines 1319--1327). The article's complete original proof of it is delivered with it (source0.tex lines 1328--1495, one \texttt{proof} environment of 168 lines). No mathematics is written, repaired or re-derived: the statement and the proof are the article's own lines, in the article's own notation, and no retained line is substituted or re-flowed.  Retained uncounted as the source-local context the proof needs: lines 640--643; lines 2416--2433.  Nothing was omitted from any retained span, so the package carries no omission record.  No retained line contains a citation, so the wrapper carries no bibliography and no bibliography entry.  No retained line contains a figure environment, an image-inclusion command or drawing code, so no image is delivered and the package declares no asset dependency.  Editorial formatting only, in addition: an AMS \texttt{amsart} preamble that restates the article's own declarations and macros used by the retained lines, namely the environments \texttt{prop}, \texttt{lem} on separate counters, together with the standard packages those lines need.  The retained environments print in this document as Proposition 1, Lemma 1 in order of appearance, whereas the article prints the same items with the numbering of its own full document.  The complete native source of the article as distributed in the arXiv e-print for this exact version is bundled unchanged as \texttt{sources/155571/source0.tex}.
\begin{equation}
\label{eq:InteractionProduct}
P_{\sigma_u \sigma_v} Q_{uv}
\end{equation}

\begin{lem}
\label{lem:conc_poisson_new}
Let $X$ be Poisson-distributed with mean $\mu \log n$. Then for any $0 < \delta < \mu$,
\[
 \P( X \le \delta \log n )
 \wle n^{- \mu h(\de/\mu)},
\]
where
$h(x) = x \log x + 1 - x$.
Furthermore, if $0 < \al < \mu-1$, then
%Furthermore, if $\mu > 1$, then for any $0<\eps'<\mu-1$ there exists a $\delta' > 0$ such that
\[
 \P( X \le \be \log n )
 \wle n^{-(1 + \al)}
\]
with $\be = \mu h^{-1}( \frac{1+\al}{\mu} )$,
where $h^{-1}(\cdot)$ denotes the inverse of $h(\cdot)$ on $(0,1)$.
\end{lem}

\begin{prop}
\label{prop:origin_bad}
For all $\la >0$, $p,q \in(0,1)$, and geometric kernel $\phi$
with a bounded normalised interaction range $\zeronorm{\phi}$ satisfying $\la \zeronorm{\phi} >1$,  if $(V^{(n)}, \bsigma^{(n)}, \bA^{(n)}) \sim \GKBM_n(\la, \phi, p, q)$ and $\la I_{\phi}(p,q) <1$, then
\[
\lim_{n \to \infty}  n \P^0(\cE_0)
\weq \infty.
\]
\end{prop}

\begin{proof}
Conditioning on the community of the node at the origin,
\begin{equation}
\P^0(\cE_0)
= \sum_{k \in \{\pm 1\}} \frac12 \P^0\bigg(\frac{\cL_0(-k)}{\cL_0(k)} \ge 1\ \Big|\  \sigma_0 = k\bigg).
\label{eq:origin_bad}
\end{equation}
Consider the term with $\sigma_0 = +1$. Then
\[
\bigg\{\frac{\cL_0(-1)}{\cL_0(1)} \ge 1 \bigg\} \weq \bigg\{\sum_{v\in V\backslash \{0 \}} \log \frac{\Ber(P_{-1, \sigma_v}Q_{0v})(A_{0v})}{\Ber(P_{1,\sigma_v}Q_{0v})(A_{0v})} \ge 0\bigg\}
\]
Since the edges $A_{0v}$ are generated with the node at the origin being in the $+1$ community, for nodes with $\sigma_v=-1$, the log-likelihood ratio evaluates to
\begin{align*}
\log \frac{\Ber(P_{-1, \sigma_v}Q_{0v})(A_{0v})}{\Ber(P_{1,\sigma_v}Q_{0v})(A_{0v})} 
&= 
\log \bigg(\Big(\frac{pQ_{0v}}{qQ_{0v}}\Big)^{A_{0v}}\Big(\frac{1-pQ_{0v}}{1-qQ_{0v}}\Big)^{1-A_{0v}}\bigg)\\
&=
A_{0v}\log \Big(\frac{pQ_{0v}}{qQ_{0v}}\Big)+ \big( 1-A_{0v} \big)\log\bigg(\frac{1-pQ_{0v}}{1-qQ_{0v}}\bigg).
\end{align*}
A similar expression with a negative sign is obtained when $\sigma_v=+1$. Combining the two, we obtain
\begin{align}
\P^0\Big( \frac{\cL_0(-1)}{\cL_0(1)} \ge 1\Big|\sigma_0=+1\Big)& \weq \P^0\Big( \sum_{v\in \cV(0)}\sigma_v\Big[A_{0v}\log\frac{q}{p}+(1-A_{0v}) \log \frac{1-qQ_{0v}}{1-pQ_{0v}}\Big]\ge 0 \,\Big|\,\sigma_0=+1\Big) \nonumber\\
&\weq \sum_{m=0}^\infty \P(|\cV(0)|=m) \ \P^0\Big( \sum_{v=1}^m R_v \ge 0 \,\Big|\,\sigma_0=+1\Big)%\sum_{v=1}^m\sigma_v\Big[A_{0v}\log\frac{q}{p}+(1-A_{0v}) \log \frac{1-qQ_{0v}}{1-pQ_{0v}}\Big]\ge 0 \ \Big| \ |\cV(0)|=m\Big).
\label{eq:map_Poisson_sum}
\end{align}
The same expression is obtained when $\sigma_0=-1$. Note that $Q_{0v} = \phi\Big(\frac{\|v\|}{\log n}\Big)$. In the following, we obtain a large deviation bound for the second probability term on the RHS. %Define $R_v^+:=\sigma_v\Big[A_{0v}\log\frac{q}{p}+(1-A_{0v}) \log \frac{1-qQ_{0v}}{1-pQ_{0v}}\Big]$ where the superscript $+$ indicates that the origin is in the $+1$ community. Our interest is in obtaining a large deviation principle (LDP) for $\P^0\big(\sum_{v=1}^m R_v^+\ge 0\big)$. 
We proceed by first computing the moment generating function of $R_v$. To indicate the conditioning event $\sigma_0=+1$, we use the notation $\P^0_+, \E^0_+$ for the conditional (Palm) probability and expectation.

Let $\zeronorm{\phi} = \ka$. Note that given the number of points in the visible set of the origin, each node is distributed uniformly within $[-\ka \log n,\ka \log n]$, assigned community $\{+1,-1\}$ independently with equal probability, and an edge is drawn to the origin based on its community and location as in \eqref{eq:InteractionProduct}. Since the same procedure is performed independently for each of
the $m$ nodes, each of the $R_v$ variables has the same distribution. Moreover, $\big\{R_v, v=1,\cdots,m\big\}$ are all independent. Integrating out the community and location of node $v$, we have that
\[ 
\E^0_+\Big[\exp(t R_v)\Big] = \frac{1}{4\ka \log n} \int_{-\ka \log n}^{\ka \log n}  \Big[\E^0_+\big[\exp(t R_v)\big| \sigma_v=-1 ,v\big]+ \E^0_+\big[\exp(t R_v)\big| \sigma_v=+1,v\big]\Big]dv.
\]
Recall that $Q_{0v} = \phi\Big(\frac{\|v\|}{\log n}\Big)$. Since $\sigma_0 = +1$, the first expectation evaluates to 
\begin{align*}
\E^0_+\big[\exp(t R_v)\big| \sigma_v=-1,v\big]
&= \E^0_+\bigg[\exp\bigg(t\sigma_v\Big[A_{0v}\log\frac{q}{p}+(1-A_{0v}) \log \frac{1-qQ_{0v}}{1-pQ_{0v}}\Big] \bigg)\bigg| \sigma_v=-1,v \bigg]\\
&= \E^0_+\bigg[\exp\bigg(t\Big[A_{0v}\log\frac{p}{q}+(1-A_{0v}) \log \frac{1-pQ_{0v}}{1-qQ_{0v}}\Big] \bigg) \bigg]\\
&=(pQ_{0v})^t (qQ_{0v})^{1-t}+(1-pQ_{0v})^t(1-qQ_{0v})^{1-t}
\end{align*}
and similarly 
\[ \E^0_+\big[\exp(t R_v)\big| \sigma_v=+1,v\big] = (qQ_{0v})^t (pQ_{0v})^{1-t}+(1-qQ_{0v})^t(1-pQ_{0v})^{1-t}.\]
Therefore, we obtain
\begin{align}
\E^0_+\Big[\exp(t R_v)\Big]
&=\frac{1}{4\ka \log n} \int_{-\ka \log n}^{\ka \log n} \Big[(pQ_{0v})^t (qQ_{0v})^{1-t}+(1-pQ_{0v})^t(1-qQ_{0v})^{1-t} \nonumber\\ 
&\hspace{3.5cm}+(qQ_{0v})^t (pQ_{0v})^{1-t}+(1-qQ_{0v})^t(1-pQ_{0v})^{1-t}\Big] dv \nonumber\\
&=\frac{1}{2\ka \log n} \int_{0}^{\ka \log n} \Big[(pQ_{0v})^t (qQ_{0v})^{1-t}+(1-pQ_{0v})^t(1-qQ_{0v})^{1-t}\nonumber\\ 
&\hspace{3.5cm}+(qQ_{0v})^t (pQ_{0v})^{1-t}+(1-qQ_{0v})^t(1-pQ_{0v})^{1-t}\Big] d\|v\|
\label{eq:mgf1}
\end{align}
Putting $\frac{\|v\|}{\log n} = z$, we get
\begin{align}
\E^0_+\Big[\exp(t R_v)\Big]
&= \frac{1}{2\ka} \int_0^{\ka } \Big[(p\phi(z))^t (q\phi(z))^{1-t}+(1-p\phi(z))^t(1-q\phi(z))^{1-t} \nonumber\\
&\hspace{3.5cm}+(q\phi(z))^t (p\phi(z))^{1-t}+(1-q\phi(z))^t(1-p\phi(z))^{1-t}\Big] dz.
\label{eq:mgf2}
\end{align}
Since the above expression is symmetric with respect to $p,q$ and $t,1-t$, the integrand is symmetric around $t=\frac{1}{2}$. Thus, the moment generating function is symmetric around $\frac{1}{2}$ within $t\in[0,1]$. Since the moment generating function is convex in $t$, it is minimized when $t=\frac{1}{2}$. Therefore, the cumulant generating function defined as $\La(t) = \log \E^0_+\Big[\exp(t R_v)\Big]$ is minimized at $t=\frac{1}{2}$. The minimum value equals 
\begin{equation}
\La\Big(\frac{1}{2}\Big)= \log \Big(\frac{1}{\ka} \int_0^\ka \Big[\sqrt{pq}\phi(z)+\sqrt{(1-p\phi(z))(1-q\phi(z))}\Big] dz\Big).
\label{eq:min_mgf}
\end{equation}
From Cram\'{e}r's theorem, for $\alpha >\E ^0_+[R_v]$,
\[
\lim_{m \to \infty} \frac{1}{m} \log \P^0_+\Big(\sum_{v=1}^m R_v \ge \alpha m \Big)
= -\La^*(\alpha).
\]
where $\La^*(\alpha)$ is the Fenchel--Legendre transform of $\La(t)$ defined as 
$\La^*(\alpha) = \sup_{t \in \R}\big[t\alpha-\La(t)\big]$.
For $\alpha = 0$, this evaluates to $\La^*(0) = -\inf_{t\in\R} \La(t)$.
Note that, using a similar procedure as in \eqref{eq:mgf1} and \eqref{eq:mgf2},
the expected value of $R_v$ can be evaluated to be
\begin{align*}
 \E^0_+[R_v] 
 \weq \frac{1}{4\ka \log n} \int_{-\ka \log n}^{\ka \log n}
   \bigg[pQ_{0v} \log \frac{qQ_{0v}}{pQ_{0v}}+(1-pQ_{0v})\log \frac{1-qQ_{0v}}{1-pQ_{0v}}\bigg]dv
 \ < \ 0,
%&= \frac{-1}{4\ka \log n} \int_{-\ka \log n}^{\ka \log n} D\bigg(\Ber\bigg(p\phi\bigg(\frac{\|v\|}{\log n}\bigg)\bigg) \ \bigg|\bigg|\  \Ber\bigg(q\phi\bigg(\frac{\|v\|}{\log n}\bigg)\bigg)\bigg) dv\\
%&<0,
\end{align*}
since the integrand is the negative of the KL divergence between two Bernoulli distributions with parameter $pQ_{0v}$ and $qQ_{0v}$.
Thus Cram\'{e}r's theorem is applicable with $\alpha = 0$.
Since $\La(\cdot)$ is convex, the infimum of the cumulant generating function $\La(t)$ 
is achieved at $t=\frac{1}{2}$ and we obtain for any $\ga>0$ and a large enough $m$
\[
\bigg| \frac{1}{m} \log\P^0_+\Big(\sum_{v=1}^m R_v \ge 0 \Big) - \La \Big(\frac{1}{2}\Big)\bigg|
\wle \ga.
% \quad 
% \text{ or }
% \quad 
% \P^0_+\big(\sum_{v=1}^m R_v \ge 0 \big) \wge \exp\big(m(\La(\frac{1}{2})-\eps)\big).
\]
A similar large deviation bound is obtained with $\P_+^0(\cdot)$ replaced by $\P_-^0(\cdot)$. %for $R_v$ when the origin is in community \(-1\). 
Using the above equation in \eqref{eq:map_Poisson_sum}, for any $\ga>0$ there exists an $m_0$ large enough such that ,
\[
 \P^0\Big( \frac{\cL_0(-1)}{\cL_0(1)}
 \wge 1\Big|\sigma_0=+1\Big)
 \wge \sum_{m=m_0}^\infty \P(|\cV(0)|=m) \exp\Big(m\Big(\La\Big(\frac{1}{2}\Big)-\ga\Big)\Big). 
\]
Including the initial terms of the summation, we obtain
\begin{align*}
 \P^0\Big( \frac{\cL_0(-1)}{\cL_0(1)} 
 &\wge 1 \Big|\sigma_0=+1\Big)\\
 &\wge  \sum_{m=0}^\infty \P(|\cV(0)|=m)\exp\Big(m\Big(\La\Big(\frac{1}{2}\Big)-\ga\Big)\Big) - \sum_{m=0}^{m_0} \P(|\cV(0)|=m)\exp\Big(m\Big(\La\Big(\frac{1}{2}\Big)-\ga\Big)\Big)\\
 &\wge \sum_{m=0}^\infty \P(|\cV(0)|=m)
  \exp\Big(m\Big(\La\Big(\frac{1}{2}\Big)-\ga\Big)\Big) - \P(|\cV(0)| \le m_0).
\end{align*}
Since $|\cV(0)|$ is a Poisson random variable with mean $2 \la \ka \log n$,
the first term is its moment generating function evaluated at
$\La(\frac{1}{2})-\ga$.
For a random variable $X\sim \Poi(\mu)$,
$\E[e^{tX}] = \exp\big(\mu (e^t-1)\big)$. 
This yields
\begin{align*}
 \P^0\Big( \frac{\cL_0(-1)}{\cL_0(1)} \ge 1\Big|\sigma_0=+1\Big)
 &\wge \exp\Big[2 \la \ka \log n \big(e^{\La(\frac{1}{2})-\ga}-1\big) \Big]-\P(|\cV(0)| \le m_0)\\ 
 &\weq n^{2 \la \ka \big(e^{\La(\frac{1}{2})-\ga}-1\big)} -\P(|\cV(0)| \le m_0).
\end{align*}
A similar computation results in 
\[
 \P^0\Big( \frac{\cL_0(+1)}{\cL_0(-1)} \ge 1\Big|\sigma_0=-1\Big)
 \wge n^{2 \la \ka \big(e^{\La(\frac{1}{2})-\ga}-1\big)} -\P(|\cV(0)| \le m_0),
\]
which together when substituted in \eqref{eq:origin_bad} yields 
\begin{equation}
 \label{eq:e0withm0}
 \P^0(\cE_0 )
 \wge n^{2 \la \ka \big(e^{\La(\frac{1}{2})-\ga}-1\big)} -\P(|\cV(0)| \le m_0),
\end{equation}
Since $e^{\La(\frac{1}{2})} =1-\frac{I_{\phi}(p,q)}{2\ka}$ from \eqref{eq:min_mgf}, and 
\[
\lim_{\ga \to 0} 2 \la \ka \big(e^{\La(\frac{1}{2})-\ga}-1\big) \weq 2 \la \ka \big(e^{\La(\frac{1}{2})}-1\big) \weq -\la I_{\phi}(p,q)>-1,
\]
taking $\beta = \frac{1+2 \la \ka \big(e^{\La(\frac{1}{2})}-1\big)}{3} > 0$, we can choose a $\ga$ small enough such that 
\begin{equation}
2 \la \ka \big(e^{\La(\frac{1}{2})-\ga}-1\big) > -1+2\beta.
\label{eq:beta_term}
\end{equation}
Using \eqref{eq:beta_term} in \eqref{eq:e0withm0}, we obtain 
\begin{equation}
\P^0(\cE_0 )
\wge n^{-1+2\beta} -\P(|\cV(0)| \le m_0).
\label{eq:e0withbetam0}
\end{equation}
The latter term in \eqref{eq:e0withbetam0} is the tail probability of a Poisson random variable whose mean is $2\la \ka \log n$. Let $c<2\beta$ be a constant and $m_0 = \ga' \log n$. We will show that for an appropriate choice of $\ga'$, $\P(|\cV(0)| \le \ga' \log n) \le n^{-1+c}$. Indeed,
using a standard Chernoff bound (Lemma~\ref{lem:conc_poisson_new}), we obtain
\[
\P(|\cV(0)| \le m_0) 
\weq \P(|\cV(0)| \le \ga' \log n)
\wle n^{-2 \la \ka h\big(\frac{\ga'}{2 \la \ka}\big)},
\]
where
$h(x) = x \log x + 1 - x$.
% \begin{align*}
% \P(|\cV(0)| \le m_0) 
% &=\P(|\cV(0)| \le \ga \log n) \\
% &\le \exp\bigg[-\ga \log n \log \Big(\frac{\ga}{2\la \ka }\Big)-(2\la \ka-\ga)\log n\Big)\bigg]\\
% &\le n^{-f(\ga)},
% \end{align*}
% where $f(\ga) = \ga \log (\frac{\ga}{2\la \ka })+(2\la \ka-\ga)$. 
Note that $\lim_{\ga'\to 0}h\big(\frac{\ga'}{2 \la \ka}\big) = 1$, and $h\big(\frac{\ga'}{2 \la \ka}\big)$ is strictly decreasing for $0<\ga'<2\la \ka$. Since $2\la \ka \ge 1$, for sufficiently small $\ga'$, $2 \la \ka h\big(\frac{\ga'}{2 \la \ka}\big) >1-c$ and we obtain $\P(|\cV(0)| \le m_0) \le n^{-1+c}$. Substituting in \eqref{eq:e0withbetam0}, since $c<2\beta$, we can write
\[
\P^0(\cE_0 )
\wge n^{-1+\beta},
\]
where $\beta>0$ whenever $\la I_{\phi}(p,q)<1$.
\end{proof}

\end{document}
